\section{Выводы}
Выполнив третью лабораторную работу по курсу \enquote{Информационный поиск}, я
реализовал токенизацию текстов документов --- компонент, на котором будет
строиться индекс следующих лабораторных работ. Основное, что я узнал нового:

\begin{itemize}
    \item выбор правил токенизации --- компромисс: простое разбиение по
    разделителям предсказуемо и не требует словарей, но апострофы, дефисы и
    составные идентификаторы (UUID) режутся, а список разделителей приходится
    расширять по факту на реальном корпусе;
    \item токенизация реального корпуса --- это в первую очередь препроцессинг:
    из 260 МБ сырого HTML в токены попадает 41{,}9 МБ текста, и без корректного
    вырезания \texttt{<script>/<style>}/комментариев статистика оказывается
    мусорной;
    \item UTF-8 нельзя обрабатывать как \texttt{char}: многобайтовые символы не
    помещаются в символьные литералы, кириллицу нужно приводить к нижнему
    регистру по байтам (и легко ошибиться: \texttt{ё} и \texttt{б} отличаются
    одним байтом), а юникодные пробелы (U+2000--U+205F) склеивают слова, если
    не считать их разделителями;
    \item эвристические фильтры (длина токена, \enquote{только цифры}, стоп-слова)
    дёшевы и убирают заметный шум, но каждый --- источник ложных срабатываний:
    \texttt{b-side} теряет \texttt{b}, таймстампы режутся в даты;
    \item скорость токенизации (11{,}7 МБ/с) не оптимальна и упирается не в
    сам проход по тексту, а в лишние копии строк и 17 find/replace-проходов
    защиты аббревиатур; ускорение --- таблица классов символов, нарезка токенов
    без копирования и многопоточность по файлам.
\end{itemize}

Итог: из 30000 документов корпуса выделено 5{,}5 млн токенов (183 на документ,
средняя длина 7{,}61 символа) за 7{,}6 с. Полученный поток токенов готов к
использованию при построении индекса в следующей лабораторной работе.

\pagebreak
